\chapter{ปัญญาประดิษฐ์วิทัศน์ถาดอาหารกุ้งและการควบคุมเครื่องให้อาหารอัตโนมัติ}
\label{chap:ch07}

\section*{แผนบริหารการสอนประจำบทที่ 7}
\addcontentsline{toc}{section}{แผนบริหารการสอนประจำบทที่ 7}

\begin{planbox}{หัวข้อเนื้อหาประจำบท}
\begin{enumerate}[topsep=2pt, itemsep=1pt, partopsep=0pt, parsep=0pt, leftmargin=1.5em]
    \item ความสำคัญทางเศรษฐศาสตร์ของการจัดการอาหารกุ้งและผลกระทบต่อคุณภาพน้ำ
    \item สถาปัตยกรรมระบบกล้องตรวจยอยกอาหารกุ้งอัตโนมัติเหนือน้ำ
    \item การประมวลผลภาพเพื่อตัดแยกและนับจำนวนเม็ดอาหารเหลือด้วยเทคนิคเซกเมนเทชัน
    \item การคำนวณอัตราส่วนพื้นที่เม็ดอาหารเหลือและการประเมินพฤติกรรมการกินอาหาร
    \item อัลกอริทึมวงปิดควบคุมเครื่องหว่านอาหารอัตโนมัติแบบปรับเปลี่ยนตามความต้องการจริง
\end{enumerate}
\end{planbox}

\begin{planbox}{วัตถุประสงค์เชิงพฤติกรรม}
\begin{enumerate}[topsep=2pt, itemsep=1pt, partopsep=0pt, parsep=0pt, leftmargin=1.5em]
    \item อธิบายความสัมพันธ์ระหว่างอาหารเหลือกับปัญหาสารประกอบไนโตรเจนที่เป็นพิษในบ่อกุ้งได้
    \item ประยุกต์ใช้วิธี Adaptive Thresholding และ Otsu ในการตัดแยกเม็ดอาหารบนตาข่ายยอได้
    \item คำนวณอัตราส่วนพื้นที่อาหารเหลือบนถาดอาหารกุ้งได้อย่างแม่นยำ
    \item พัฒนาโปรแกรมควบคุมรอบการทำงานของเครื่องหว่านอาหารแบบวงปิดได้อย่างถูกต้อง
\end{enumerate}
\end{planbox}

\newpage

\begin{planbox}{วิธีสอนและกิจกรรมการเรียนการสอน}
\begin{enumerate}[topsep=2pt, itemsep=1pt, partopsep=0pt, parsep=0pt, leftmargin=1.5em]
    \item บรรยายหลักการโภชนาการสัตว์น้ำและระบบคอมพิวเตอร์วิทัศน์ในการตรวจวัดวัตถุขนาดเล็ก
    \item สาธิตการถ่ายภาพยอยกอาหารกุ้งภายใต้สภาพแสงสะท้อนของผิวน้ำ
    \item ปฏิบัติการเขียนโค้ดภาษา Python สกัดพื้นที่เม็ดอาหารและกำจัดสัญญาณรบกวน
    \item จำลองการควบคุมเครื่องหว่านอาหารอัจฉริยะ (Smart Autofeeder\index{เครื่องให้อาหารอัตโนมัติอัจฉริยะ (Smart Autofeeder)}\index{Smart Autofeeder}) ด้วยสัญญาณ PWM
\end{enumerate}
\end{planbox}

\begin{planbox}{สื่อการเรียนการสอน}
\begin{enumerate}[topsep=2pt, itemsep=1pt, partopsep=0pt, parsep=0pt, leftmargin=1.5em]
    \item เอกสารคำสอนบทที่ 7 สไลด์บรรยาย และชุดภาพถ่ายยอยกอาหารกุ้งจริง
    \item ถาดอาหารกุ้ง (ยอยก) พร้อมเม็ดอาหารกุ้งสำเร็จรูปเบอร์ 1 ถึง 4
    \item กล้องประมวลผลภาพเอดจ์พร้อมไฟส่องสว่างโพลาไรซ์ตัดแสงสะท้อน
    \item บอร์ด ESP32-S3 ขับมอเตอร์จ่ายอาหารกุ้งและโซลินอยด์ปล่อยอาหาร
\end{enumerate}
\end{planbox}

\begin{planbox}{การวัดและประเมินผล}
\begin{enumerate}[topsep=2pt, itemsep=1pt, partopsep=0pt, parsep=0pt, leftmargin=1.5em]
    \item การทดสอบความรู้เรื่องอัลกอริทึมเซกเมนเทชันและการวิเคราะห์สัญฐานวิทยาของภาพ
    \item การประเมินความแม่นยำในการตรวจนับพื้นที่เม็ดอาหารเปรียบเทียบกับค่าจริง
    \item การทดสอบการทำงานของระบบวงปิดในการปรับเพิ่ม-ลดปริมาณอาหารอัตโนมัติ
\end{enumerate}
\end{planbox}

\clearpage

\section{ความสำคัญของการจัดการอาหารกุ้งต่อสิ่งแวดล้อมบ่อเลี้ยง}

ในต้นทุนการผลิตกุ้งขาวแวนนาไม ค่าอาหารสำเร็จรูปคิดเป็นสัดส่วนสูงถึง \textbf{ร้อยละ 60 ถึง 70} ของต้นทุนการผลิตทั้งหมด การให้อาหารอย่างมีประสิทธิภาพจึงเป็นปัจจัยชี้ขาดความอยู่รอดทางธุรกิจ

การให้อาหารมากเกินความต้องการ (Overfeeding) นอกจากจะเป็นการสิ้นเปลืองต้นทุนแล้ว อาหารที่เหลือจะตกลงสู่ก้นบ่อ เกิดการย่อยสลายโดยจุลินทรีย์แบบไม่ใช้ออกซิเจน ก่อให้เกิดก๊าซพิษ เช่น แอมโมเนียอิสระ ($\text{NH}_3$) ไนไตรท์ ($\text{NO}_2^-$) และไฮโดรเจนซัลไฟด์ ($\text{H}_2\text{S}$) ซึ่งทำลายเหงือกกุ้งและกระตุ้นการระบาดของโรคขี้ขาว

ในทางกลับกัน การให้อาหารน้อยเกินไป (Underfeeding) จะส่งผลให้อัตราการเจริญเติบโตชะงัก กุ้งมีขนาดแตกไซส์ และเกิดพฤติกรรมการกินกันเอง (Cannibalism)

\section{ระบบวิทัศน์ตรวจถาดอาหารกุ้งอัตโนมัติ}

การตรวจสอบการกินอาหารของกุ้งในปัจจุบันยังใช้วิธีแรงงานคนพายเรือไปยกยอตรวจดูอาหารเหลือทุก $2 - 2.5\text{ ชั่วโมง}$ ซึ่งมีความแปรปรวนสูงตามดุลยพินิจของพนักงาน การนำระบบรอกไฟฟ้าดึงยอยกขึ้นสู่ผิวน้ำพร้อมติดตั้งกล้องคอมพิวเตอร์วิทัศน์ที่ปลายเสา ทำให้สามารถตรวจวัดปริมาณอาหารเหลือได้แบบอัตโนมัติและแม่นยำ

\subsection{การตัดแยกเม็ดอาหารด้วยอัลกอริทึมเซกเมนเทชัน}

ภาพถ่ายยอยกอาหารกุ้งจะประกอบด้วย 3 องค์ประกอบหลัก ได้แก่ พื้นหลังตาข่ายยอ น้ำขังสะท้อนแสง และเม็ดอาหารสำเร็จรูป การตัดแยกเม็ดอาหารอาศัย \textbf{การแบ่งส่วนภาพแบบปรับตัว (Adaptive Thresholding)\index{การแบ่งส่วนภาพแบบปรับตัว (Adaptive Thresholding)}\index{Adaptive Thresholding}} ร่วมกับการคัดกรองพื้นที่สัญฐานวิทยา (Morphological Area Filtering)

\begin{equation}
T(x, y) = \text{mean}(I_{\text{local}}) - C_0
\label{eq:adaptive_thresh}
\end{equation}

โดยที่ $T(x, y)$ แทนค่าขีดเริ่มเปลี่ยนเฉพาะจุด, $I_{\text{local}}$ แทนความสว่างในหน้าต่างรอบจุดพิกเซล ($21 \times 21$ พิกเซล), และ $C_0$ แทนค่าคงที่ชดเชยความต่างสี

\begin{pythonbox}[การคำนวณอัตราส่วนพื้นที่อาหารเหลือบนถาดอาหารกุ้งด้วย Python]
import cv2
import numpy as np

def analyze_feeding_tray_pellets(image_path: str):
    image = cv2.imread(image_path)
    gray = cv2.cvtColor(image, cv2.COLOR_BGR2GRAY)
    
    # เบลอเพื่อลดสัญญาณรบกวนของตาข่ายยอ
    blurred = cv2.GaussianBlur(gray, (5, 5), 0)
    
    # ใช้วิธี Adaptive Thresholding ตัดแยกเม็ดอาหาร
    thresh = cv2.adaptiveThreshold(
        blurred, 255, cv2.ADAPTIVE_THRESH_GAUSSIAN_C,
        cv2.THRESH_BINARY_INV, 21, 4
    )
    
    # ตรวจหาเส้นขอบของเม็ดอาหารและกรองขนาดวัตถุ
    contours, _ = cv2.findContours(thresh, cv2.RETR_EXTERNAL, cv2.CHAIN_APPROX_SIMPLE)
    
    total_pellet_area = 0
    pellet_count = 0
    # กำหนดพื้นที่เม็ดอาหารกุ้งจริง (ระหว่าง 15 ถึง 250 พิกเซล)
    for c in contours:
        area = cv2.contourArea(c)
        if 15 <= area <= 250:
            total_pellet_area += area
            pellet_count += 1
            
    # คำนวณอัตราส่วนพื้นที่อาหารเทียบกับพื้นที่ถาดทั้งหมด
    tray_area = gray.shape[0] * gray.shape[1]
    uneaten_ratio = total_pellet_area / tray_area
    
    return {
        "pellet_count": pellet_count,
        "uneaten_ratio": uneaten_ratio,
        "pellet_area_px": total_pellet_area
    }
\end{pythonbox}

\section{การควบคุมเครื่องหว่านอาหารอัตโนมัติแบบวงปิด}

ข้อมูลอัตราส่วนอาหารเหลือ ($R_{\text{uneaten}}$) จะถูกส่งต่อไปยังคอนโทรลเลอร์เพื่อปรับเวลาการเปิดวาล์วหว่านอาหารของ \textbf{เครื่องให้อาหารอัตโนมัติ (Smart Autofeeder)} ในรอบถัดไปแบบวงปิด (Closed-Loop Feeding Control)

\begin{table}[H]
\centering
\caption{กฎเกณฑ์การปรับเปลี่ยนปริมาณอาหารของระบบ Smart Autofeeder}
\label{tab:autofeeder_rules}
\begin{tabularx}{\textwidth}{p{3.2cm}p{4.5cm}Y}
    \toprule
    \textbf{อัตราส่วนอาหารเหลือ ($R_{\text{uneaten}}$)} & \textbf{สภาวะการกินอาหาร} & \textbf{คำสั่งควบคุมเครื่องให้อาหารในรอบถัดไป} \\
    \midrule
    $< 0.03$ (อาหารเกลี้ยงเร็ว) & กุ้งกินอาหารดีมาก อาหารไม่เพียงพอ & เพิ่มปริมาณอาหารขึ้น $+15\%$ (เพิ่มระยะเวลาหว่าน) \\
    $0.03 - 0.08$ & กุ้งกินอาหารในเกณฑ์มาตรฐาน & คงปริมาณอาหารเดิมตามตารางคำนวณน้ำหนักตัว \\
    $0.08 - 0.15$ & เริ่มมีอาหารเหลือตกค้าง & ปรับลดปริมาณอาหารลง $-20\%$ \\
    $> 0.15$ (อาหารเหลือมาก) & กุ้งหยุดกิน อาจเกิดปัญหาน้ำเสียหรือโรค & หยุดการหว่านอาหารทันที และส่งสัญญาณแจ้งเตือน \\
    \bottomrule
\end{tabularx}
\end{table}

\begin{theorybox}[การประสานเวลาการให้อาหารกับระดับออกซิเจนละลายน้ำ]
ระบบ Smart Autofeeder จะทำงานร่วมกับหัววัด Optical DO เสมอ:
\begin{itemize}[leftmargin=1.5em]
    \item หากค่า DO ต่ำกว่า $4.0\text{ mg/L}$ ระบบจะระงับการจ่ายอาหารโดยอัตโนมัติ แม้ว่าถาดอาหารจะไม่มีอาหารเหลือก็ตาม เนื่องจากกุ้งที่กินอาหารในสภาวะออกซิเจนต่ำจะเกิดภาวะเครียดและย่อยอาหารไม่ได้
    \item การประสานงานร่วมกันนี้ช่วยลดการสูญเสียอาหารลงได้ร้อยละ 18 ถึง 25 ต่อรอบการเลี้ยง และลดค่าแอมโมเนียสะสมในบ่อลงได้อย่างมีนัยสำคัญ
\end{itemize}
\end{theorybox}

\section*{คำถามทบทวนท้ายบทที่ 7}
\addcontentsline{toc}{section}{คำถามทบทวนท้ายบทที่ 7}

\begin{enumerate}
    \item การให้อาหารกุ้งเกินความต้องการส่งผลกระทบต่อวัฏจักรไนโตรเจนและคุณภาพน้ำในบ่อเลี้ยงอย่างไร
    \item จงอธิบายหลักการของ Adaptive Thresholding และเปรียบเทียบข้อดีกับ Global Thresholding (Otsu) ในกรณีภาพถ่ายยอยกอาหารที่มีแสงสะท้อนไม่สม่ำเสมอ
    \item หากตรวจพบอัตราส่วนอาหารเหลือบนถาดเท่ากับ 0.02 ในขณะที่ค่า DO ในบ่อเท่ากับ 3.6 mg/L ระบบอัจฉริยะควรตัดสินใจสั่งการอย่างไร พร้อมให้เหตุผลทางสรีรวิทยา
    \item จงคำนวณมูลค่าการประหยัดต้นทุนในฟาร์มกุ้ง 10 บ่อ ที่มีต้นทุนอาหารรวม 2,000,000 บาทต่อรอบ หากระบบ AIoT สามารถลดการสูญเสียอาหารลงได้ร้อยละ 20
\end{enumerate}
\clearpage
